% cover_letter.tex — Cover letter for Wireless Personal Communications (R1 → R2 revision)
% Compile: bash build.sh cover
\documentclass[11pt,a4paper]{article}
\usepackage[margin=2.5cm]{geometry}
\usepackage[utf8]{inputenc}
\usepackage[T1]{fontenc}
\usepackage{lmodern}
\usepackage{amsmath,amssymb}
\usepackage[mode=text]{siunitx}
\sisetup{locale=US}
\usepackage{hyperref}
\hypersetup{hidelinks}

\setlength{\parindent}{0pt}
\setlength{\parskip}{0.6em}

\begin{document}

\begin{center}
{\Large\bfseries Cover Letter}\\[4pt]
{\itshape Revised submission to \emph{Wireless Personal Communications} (R1 $\to$ R2)}
\end{center}

\vspace{1em}

\noindent\textbf{To:} The Editor-in-Chief, \emph{Wireless Personal Communications} (Springer)\\
\textbf{From:} Bin Nong, Weihong Fu, and Zhuoyun Jiang\\
\textbf{Corresponding email:} \href{mailto:nongbin@tfswufe.edu.cn}{nongbin@tfswufe.edu.cn}\\
\textbf{Manuscript title:} A Lightweight Complex-Valued CNN with Complex Squeeze-and-Excitation Attention for Single-Channel Blind Source Separation of Co-Frequency Communication Signals\\
\textbf{Manuscript type:} Original Research Paper\\
\textbf{Revision:} R1 $\to$ R2

\vspace{1em}

\noindent Dear Editor,

\vspace{0.4em}

We thank the editor and the reviewers for the thorough and constructive critique of our original submission. The R1 review identified genuine weaknesses in the manuscript --- the missing multipath in Eq.~(1), the under-tested ``co-frequency'' claim, the soft baseline set, the insufficient statistical reporting, and several overclaims --- and we have addressed every one of them in this revision. A detailed point-by-point response is provided in \texttt{reviewer\_response.md} (27 items, R1.1--R1.12, R2.1--R2.10, R3.1--R3.3). Below we summarise the substantive changes for the editor's convenience.

\paragraph{Major changes since R1.}
\begin{enumerate}
\item \textbf{Expanded baselines.} The R1 manuscript compared only against the no-SE ablation, a real-valued baseline, and a complex Conv-TasNet. In R2 we add (i) two \emph{parameter-matched} complex-CNN baselines (no-SE at $H{=}70$ giving 242K params, real-valued at $H{=}80$, $L{=}12$ giving 237K params) to disentangle the contribution of the SE block from the parameter budget, and (ii) two \emph{cross-domain} communication-signal baselines (CNSE by Hou \& Gao 2022 and S4-UNET by Gao et al.~2026) re-implemented from the primary sources and scaled down to fit the \SI{8}{GB} GPU used in our setting.
\item \textbf{Statistical reporting.} Every evaluable configuration is now reported over \textbf{five} random seeds (42--46) with sample standard deviation. We add a dedicated subsection (Section 4.8) reporting paired Wilcoxon signed-rank tests between C-SE and every baseline on per-seed SDR. We are explicit that \emph{no comparison reaches statistical significance at $\alpha=0.05$} --- the smallest attainable two-sided $p$ at $N{=}5$ pairs is $0.0625$ --- and reposition the paper's narrative around per-parameter efficiency and per-seed reproducibility rather than raw SDR dominance.
\item \textbf{Equation (1) and scope.} Eq.~(1) now includes the multipath channel $h_k$ with explicit 3-tap complex Gaussian definition, normalisation, and unit-energy constraint. A new paragraph after the equation states the scope: the generator uses symbol-aligned mixtures with per-source carrier offset within $\pm 5$ Hz; relaxing symbol alignment is left to future work.
\item \textbf{Generalisation to the co-frequency limit.} We trained two new models (proposed C-SE and matched no-SE) with the carrier-frequency gap drawn from $U(0, 5)$ Hz per sample. Evaluating at gaps from $0$ to $500$ Hz (new Section 4.7, Fig.~micro\_freq) shows the proposed C-SE holds $\text{SDR}\approx\SI{3.3}{dB}$ across the entire range, including the previously unseen co-frequency limit $\Delta f = \SI{0}{Hz}$ ($\text{SDR}=\SI{3.33}{dB}$). The matched no-SE baseline trained under the same schedule still collapses to $\text{SIR}\approx\SI{21}{dB}$ at every gap, confirming the gap is structural.
\item \textbf{Honest positioning.} ``First'' is now softened throughout to ``to the best of our knowledge, the first phase-preserving, real-gated complex-valued SE design specifically evaluated for lightweight single-channel blind separation of near-co-frequency communication signals.'' ``SOTA'' is removed from the Conv-TasNet label. ``Is in fact optimal'' is replaced with ``is sufficient'' for the scale-mode ablation. The Discussion section acknowledges prior complex-domain work (CSENet 2022, complex-valued SE 2024/2025) and identifies the C-SE's specific contribution as the phase-preserving real-gated design.
\end{enumerate}

\paragraph{Headline numbers (5 seeds, $\text{SNR}=\SI{10}{dB}$).}
The proposed C-SE reaches $\text{SDR}=\SI{2.31 \pm 0.62}{dB}$ at \qty{235}{K} parameters (5 seeds). The matched no-SE complex CNN reaches $\SI{1.81 \pm 0.48}{dB}$ (5 seeds); the matched real-valued CNN reaches $\SI{2.20 \pm 0.59}{dB}$ (5 seeds); the Complex Conv-TasNet reaches $\SI{3.23 \pm 0.40}{dB}$ (3 seeds, 817K params). CNSE (6.69M) reaches $\SI{2.38 \pm 0.91}{dB}$ (3 seeds) and S4-UNET (1.57M) reaches $\SI{3.05 \pm 0.00}{dB}$ (3 seeds). C-SE is therefore competitive at the \qty{235}{K}-parameter scale on the dimensions of partial collapse-resistance (2 of 5 vs.\ 4 of 5 for matched no-SE), per-parameter efficiency, and predictable runtime; we make no claim of statistical dominance.

\paragraph{Reproducibility and ethics.}
\begin{itemize}
\item The manuscript reports per-seed numbers for every evaluated configuration; the full implementation will be released as open-source code (MIT) upon acceptance.
\item \texttt{docs/EXPERIMENT\_LOG.md} at the project root provides a complete reproducibility record with exact commands, hyperparameters, and per-seed SDR/SI-SDR/SIR values for all 7 experimental phases and 44 training runs.
\item \texttt{paper1\_cnn\_se/data\_radioml.py} ships an optional loader for the public RadioML 2016.10A corpus (not used in headline results; left for the community to extend the evaluation to over-the-air and public I/Q data).
\item The manuscript is original, has not been published elsewhere, and is not under consideration by any other journal.
\item The authors have no conflicts of interest to declare.
\item No human subjects or animal experiments are involved.
\end{itemize}

\paragraph{Author contributions.}
B.~Nong conceived the C-SE design, ran the multi-seed experiments, and drafted the manuscript. W.~Fu contributed the matched-baseline analysis and the scale-mode / pooling ablations. Z.~Jiang contributed the cross-domain baseline re-implementations (CNSE, S4-UNET) and the statistical-significance analysis. All authors reviewed and approved the final manuscript.

\paragraph{Funding.}
This work received no specific grant from any funding agency in the public, commercial, or not-for-profit sectors.

\paragraph{Suggested reviewers.}
We have no specific reviewer suggestions; we trust the editorial board to identify expert reviewers in signal processing for wireless communications.

\vspace{1em}

\noindent Thank you for considering our revised manuscript. We respectfully request a re-evaluation.

\vspace{2em}

\noindent Sincerely,\\[1.4em]
\textbf{Bin Nong}\\
Tianfu College of Southwestern University of Finance and Economics, Mianyang, China\\
\texttt{nongbin@tfswufe.edu.cn}\\[0.6em]
\textbf{Weihong Fu}\\
School of Telecommunications Engineering, Xidian University, Xi'an, Shaanxi, China\\[0.6em]
\textbf{Zhuoyun Jiang}\\
School of Health Science and Engineering, University of Shanghai for Science and Technology, Shanghai, China

\end{document}
